% -*- TeX:RNW:UK -*-
\documentclass[10pt]{beamer}\usepackage[]{graphicx}\usepackage[]{color}
% maxwidth is the original width if it is less than linewidth
% otherwise use linewidth (to make sure the graphics do not exceed the margin)
\makeatletter
\def\maxwidth{ %
  \ifdim\Gin@nat@width>\linewidth
    \linewidth
  \else
    \Gin@nat@width
  \fi
}
\makeatother

\definecolor{fgcolor}{rgb}{0.345, 0.345, 0.345}
\makeatletter
\@ifundefined{AddToHook}{}{\AddToHook{package/xcolor/after}{\definecolor{fgcolor}{rgb}{0.345, 0.345, 0.345}}}
\makeatother
\newcommand{\hlnum}[1]{\textcolor[rgb]{0.686,0.059,0.569}{#1}}%
\newcommand{\hlstr}[1]{\textcolor[rgb]{0.192,0.494,0.8}{#1}}%
\newcommand{\hlcom}[1]{\textcolor[rgb]{0.678,0.584,0.686}{\textit{#1}}}%
\newcommand{\hlopt}[1]{\textcolor[rgb]{0,0,0}{#1}}%
\newcommand{\hlstd}[1]{\textcolor[rgb]{0.345,0.345,0.345}{#1}}%
\newcommand{\hlkwa}[1]{\textcolor[rgb]{0.161,0.373,0.58}{\textbf{#1}}}%
\newcommand{\hlkwb}[1]{\textcolor[rgb]{0.69,0.353,0.396}{#1}}%
\newcommand{\hlkwc}[1]{\textcolor[rgb]{0.333,0.667,0.333}{#1}}%
\newcommand{\hlkwd}[1]{\textcolor[rgb]{0.737,0.353,0.396}{\textbf{#1}}}%
\let\hlipl\hlkwb

\usepackage{framed}
\makeatletter
\newenvironment{kframe}{%
 \def\at@end@of@kframe{}%
 \ifinner\ifhmode%
  \def\at@end@of@kframe{\end{minipage}}%
  \begin{minipage}{\columnwidth}%
 \fi\fi%
 \def\FrameCommand##1{\hskip\@totalleftmargin \hskip-\fboxsep
 \colorbox{shadecolor}{##1}\hskip-\fboxsep
     % There is no \\@totalrightmargin, so:
     \hskip-\linewidth \hskip-\@totalleftmargin \hskip\columnwidth}%
 \MakeFramed {\advance\hsize-\width
   \@totalleftmargin\z@ \linewidth\hsize
   \@setminipage}}%
 {\par\unskip\endMakeFramed%
 \at@end@of@kframe}
\makeatother

\definecolor{shadecolor}{rgb}{.97, .97, .97}
\definecolor{messagecolor}{rgb}{0, 0, 0}
\definecolor{warningcolor}{rgb}{1, 0, 1}
\definecolor{errorcolor}{rgb}{1, 0, 0}
\makeatletter
\@ifundefined{AddToHook}{}{\AddToHook{package/xcolor/after}{
\definecolor{shadecolor}{rgb}{.97, .97, .97}
\definecolor{messagecolor}{rgb}{0, 0, 0}
\definecolor{warningcolor}{rgb}{1, 0, 1}
\definecolor{errorcolor}{rgb}{1, 0, 0}
}}
\makeatother
\newenvironment{knitrout}{}{} % an empty environment to be redefined in TeX

\usepackage{alltt}
\usetheme{metropolis}
%\useinnertheme{rectangles}
\setbeamercovered{%
still covered={\opaqueness<1->{15}},
again covered={\opaqueness<1->{40}}}

\hypersetup{colorlinks,linkcolor=black,urlcolor=brown,citecolor=brown}

\include{shared-tex/set-up-fonts-and-icons}
\IfFileExists{upquote.sty}{\usepackage{upquote}}{}
\begin{document}





\title{IPS-141\\Sensory and Physiological Ecology of  Plants}
\subtitle{x: Stomata and gas-exchange}
\author{Pedro J. Aphalo}
\date{January--February 2022}
\institute[Univ.\ of Helsinki]{M.Sc.\ in Plant Biology, University of Helsinki\\[2ex] \url{http://blogs.helsinki.fi/aphalo/}}

  \begin{frame}
    \maketitle
  \end{frame}

  \begin{frame}[c]
    \begin{center}
      \begin{small}
        \copyright 2006--2022 by Pedro J. Aphalo\\
        University of Helsinki, Finland.\\
        \textcolor{blue}{\url{http://blogs.helsinki.fi/senpep-blog/}}\\[2ex]
      \end{small}

      \begin{footnotesize}
        Sensory and Physiological Ecology of Plants slides by Pedro J. Aphalo are licensed under a Creative Commons Attribution-ShareAlike 4.0 International License.

      \includegraphics[width=6em]{../figures/copyright/by-sa}\\[2ex]
      \end{footnotesize}
        
        \begin{scriptsize}
        Typeset in Lucida Sans, \textrm{Luicda Bright}, \texttt{Lucida Console} and Lucida Math. Icons from fonts ``WebHostingHub Glyphs'' (under SIL-Open Font License) from \url{https://www.webhostinghub.com/}; ``insect icons'' (free from \url{http://www.woodcutter.es/}); ``linea-basic-10'' and ``linea-weather-10'' (free from \url{https://github.com/linea-io}), ``Mini Pics Uprooted Twig'' and ``Mini Pics Uprooted Twig'' (commercial, from Image Club Graphics, Inc.). Plant icon as .svg by Abdul Wahhab (free from \url{NounProject.com}).

        Illustrations and text quoted from copyrighted sources is excluded from this license and their use should respect the original licenses.
        \end{scriptsize}
    \end{center}
  \end{frame}



  \begin{frame}
    \frametitle{Outline}
    \tableofcontents
  \end{frame}

\section{Stomata}

\begin{frame}{Stomata on the abaxial epidermis (\emph{Hedera helix} L.)}
 \centering
    \includegraphics[width=0.95\textwidth]{photos/ivy-leaf-abax-epidermis-lowres}
\end{frame}

\begin{frame}{Cross-section of a juvenile leaf (\emph{Hedera helix} L.)}
 \centering
    \includegraphics[width=0.95\textwidth]{photos/Ivy-leaf-section}
\end{frame}

\begin{frame}{Cross-section of a stoma (\emph{Hedera helix} L.)}
 \centering
    \includegraphics[width=0.95\textwidth]{photos/ivy-guard-cell-section}
\end{frame}

\section{Stomata, responses to light}

\begin{frame}{Time course of \gs (\emph{A. thalina})}
\begin{itemize}
  \item Experiment with \emph{phot1} \emph{phot2} and \emph{cry1} \emph{cry2} mutants plus two WTs.
  \item Growth under photoperiod 12 h, 200 \umol PAR from white light.
  \item Gas-exchange measured continuously for 23~h in whole rosettes (1 rosette at a time).
  \item Measurements in same photoperiod, 200 \umol PAR from red, green or blue light from LEDs, and in darkness.
  \item \autocite[See][for details]{Wang2020}
\end{itemize}
\end{frame}

\begin{frame}{Time course of \gs (\emph{A. thalina})}
 \centering
    \includegraphics[width=0.85\textwidth]{figures/Wang-A-gs-time-course-dark}
\end{frame}

\begin{frame}{Time course of \gs (\emph{A. thalina})}
 \centering
    \includegraphics[width=0.95\textwidth]{figures/Wang-gs-1st-hour-light-dark}
\end{frame}

\begin{frame}{Time course of \gs (\emph{A. thalina})}
 \centering
    \includegraphics[width=0.85\textwidth]{figures/Wang-A-gs-time-course-light}
\end{frame}

\begin{frame}{Time course of \gs (\emph{A. thalina})}
 \centering
    \includegraphics[width=0.95\textwidth]{figures/Wang-CiCa-time-course-light-dark}
\end{frame}

\begin{frame}{Blue component of sunlight}
 \centering
    \includegraphics[width=0.95\textwidth]{figures/blue-hour-plot}
\end{frame}

\section{Measurement of leaf conductance}

\begin{frame}{Transient porometer}
     \href{https://youtu.be/kA7yS8JJsLE}{Video describing the AP4 porometer} from \href{https://delta-t.co.uk/product/ap4/}{Delta-T}
\end{frame}

\section{Thermography and non-contact sensors}

\begin{frame}{Thermal camera}
 \centering
    \includegraphics[width=0.65\textwidth]{photos/OPTPI4XXi}
     from \href{https://www.optris.co.uk/thermal-imager-optris-pi400-pi450}{Optris}
\end{frame}

\section{Measurement of spectral irradiance}

\begin{frame}{Spectroradiometer}
     \href{https://youtu.be/OI3pIvLhVcc}{Video describing a spectrometer} from \href{https://www.oceaninsight.com/}{Ocean Insight}
\end{frame}

\section{Measurement of gas-exchange}

\begin{frame}{Open gas-exchange system}
\begin{footnotesize}
\begin{center}
\makebox[2cm]{{\sf%
air, \COtwo}%
\rightarrowfill}
\fbox{\parbox{2.3cm}{\sf%
AIR\\ CONDITIONING}}%
$\makebox[3.1cm]{\rightarrowfill} \atop
\makebox[3.1cm]{%
\rightarrowfill
\fbox{\parbox{2.1cm}{\sf%
LEAF CHAMBER}}%
\rightarrowfill}$
\fbox{\parbox{1.8cm}{\sf%
MEASURING\\ EQUIPMENT}}%
\end{center}
\end{footnotesize}
    \begin{itemize}
        \item Most current gas-exchange systems used for measuring photosynthesis
        are open systems in which air flows at a constant rate through a chamber
        where a leaf in enclosed.

        \item The flow rates and the differences in concentration of \COtwo and
        water vapour between the air entering the chamber and that coming
        out of the chamber are used to calculate the flows in or out of the leaf.

        \item From the flows of \COtwo and water in or out of the leaf, and the enclosed leaf area, the
        \COtwo and water fluxes per unit leaf area can be calculated.
    \end{itemize}
\end{frame}

\begin{frame}{Open vs.\ closed systems}
    \begin{description}
        \item[Open system:] Environmental conditions in the chamber do not
        vary during measurement.
        \item[Closed system:] Environmental conditions in the chamber change during measurement.
        With transpiration humidity increases, with photosynthesis \COtwo
        concentration decreases.
    \end{description}
\end{frame}

\begin{frame}{Units}
    \begin{itemize}
        \item Modern: using moles
            \begin{description}
                \item[\COtwo] \umol
                \item[Water vapour] \mmol
            \end{description}
        \item Old: mass
            \begin{description}
                \item[\COtwo] mg cm$^{-2}$ min$^{-1}$ etc.
                \item[Water vapour] g cm$^{-2}$ min$^{-1}$ etc.
            \end{description}
    \end{itemize}
\end{frame}

\begin{frame}{Gas-exchange system}
    \centering
    \includegraphics[height=0.7\textheight]{photos/WalzGFS-3000}\\
    {\small GFS-3000 (Walz, Germany) (\url{http://www.walz.com/}).}
\end{frame}

\begin{frame}{Measuring}
    \centering
    \includegraphics[height=0.7\textheight]{photos/GFS-3000_1big}\\
    {\small from \url{http://www.walz.com/}.}
\end{frame}

\begin{frame}{Measuring}
    \centering
    \includegraphics[height=0.7\textheight]{photos/GFS-3000_6big}\\
    {\small from \url{http://www.walz.com/}.}
\end{frame}

\begin{frame}{Chamber}
    \centering
    \includegraphics[height=0.7\textheight]{photos/GFS-3000_3big}\\
    {\small from \url{http://www.walz.com/}.}
\end{frame}

\begin{frame}{Chamber (measuring)}
    \centering
    \includegraphics[height=0.7\textheight]{photos/gFS-3000_4big}\\
    {\small from \url{http://www.walz.com/}.}
\end{frame}

\begin{frame}{Standard and conifer chambers}
    \centering
    \includegraphics[height=0.5\textheight]{photos/LI-6400-standard}%
    \includegraphics[height=0.5\textheight]{photos/LI-6400-conifer}\\
    {\small Standard chamber (for flat leaves) and conifer chamber (for twigs)
     for the LI-6400 gas-exchange system\\
     (Li-Cor, USA) (from \url{http://www.licor.com/}).}
\end{frame}


\section*{References}
  \nocite{Kotilainen2020}
  \begin{frame}[t,allowframebreaks]
    \frametitle{References}
    \printbibliography
  \end{frame}

\end{document}
